\documentclass[fontsize=9pt]{article}
\usepackage[margin=0.55in]{geometry}
\usepackage{lipsum,mwe}
\usepackage[T1]{fontenc}
\usepackage[english]{babel}
\setlength\parindent{0pt}
\usepackage{amsmath,amsfonts}
\usepackage{graphicx}
\usepackage{float}
\usepackage{subcaption}
\usepackage{enumitem}
\usepackage{xcolor}
\usepackage{eurosym}
\usepackage{hyperref}
\hypersetup{colorlinks=true, urlcolor=blue!50!black, linkcolor=black}
\usepackage{sectsty}
\allsectionsfont{\normalfont \small \scshape}
\setlength{\parskip}{2pt plus 1pt}
\pagestyle{empty}

\begin{document}

\begin{center}
{\large\textbf{NS-RAM Co-Design --- Differentiable Software Bridge from Silicon to Algorithm}} \\[2pt]
{\footnotesize One-page proposal summary.
E.~Bergvall, R.~Luciani, S.~Pazos, M.~Lanza --- \today}
\end{center}

\vspace{-2pt}
\hrule
\vspace{4pt}

\textbf{Premise.} The NS-RAM 2T floating-body cell delivers synapse,
neuron and short-term-memory behaviour in one transistor pair, in
unmodified bulk $130$~nm CMOS, at sub-pJ per cycle and tens of
square micrometres per cell. Today every algorithm-driven
cell-parameter question round-trips through hand-curated SPICE decks
and silicon characterisation, with months of latency. We propose to
fund the missing software bridge: a differentiable, GPU-batched
Python port of the BSIM4 channel and the parasitic NPN, plus a
body-state-aware network simulator, that closes the loop between
measured silicon and the next tape-out (see co-design loop figure).

\begin{figure}[H]
  \centering
  \includegraphics[width=0.78\linewidth]{figures/codesign_loop.pdf}
  \vspace{-3pt}
  \caption*{\footnotesize\emph{Co-design loop.} Measured silicon
  $\to$ differentiable simulator $\to$ next tape-out, with refit
  feedback. Without the centre box, every what-if costs a fab
  cycle.}
  \vspace{-6pt}
\end{figure}

\textbf{What we have built.}
A single Python package, GPU-batched through PyTorch on commodity
integrated graphics. Three layers: (i) a differentiable BSIM4 v4.8.3
channel against the PTM~130 card, cross-checked against ngspice;
(ii) a parasitic NPN (Gummel--Poon, with VBIC level-4 as swap-in)
tied to the channel through impact ionisation and floating-body
charging; (iii) a network backend that batches cells across
topologies (Erd\H{o}s--R\'enyi, small-world, lattice, reservoir
families) up to $N\!\geq\!10^{5}$. Throughput
$\sim\!5\!\times\!10^{6}$ cell-evaluations~s$^{-1}$ on an AMD
Radeon~$8060$S iGPU.

\begin{figure}[H]
  \centering
  \includegraphics[width=0.99\linewidth]{figures/brief_headlines_honest/brief_headlines.pdf}
  \vspace{-4pt}
  \caption*{\footnotesize\emph{Three honest headlines of the
  artifact today.} Left: cell-wide DC $\log_{10}|I_d|$~RMSE
  $\approx 1.19$~dec on the full $25$-bias set, forward$+$backward
  (no direction-pick). Centre: six of nine pre-registered dynamic
  behaviours pass on the calibrated cell. Right: seven network
  demonstrations cleared pre-registered gates on independent
  benchmarks.}
  \vspace{-6pt}
\end{figure}

\textbf{Where the artifact stands.}

\begin{enumerate}[leftmargin=1.6em,topsep=2pt,itemsep=2pt]
\item \textbf{Cell, DC.} Cell-wide $\log_{10}|I_d|$~RMSE
$\approx\!1.19$~decade on $25$ measured biases, forward$+$backward
averaged, pseudo-transient solver. This is the publishable headline;
earlier $<\!1.0$ values were forward-only or bias-subsets and have
been retracted.

\item \textbf{Cell, dynamics.} Six of nine pre-registered dynamic
tests pass: DC characteristic, forward/backward hysteresis,
nanosecond drain-snap, body-voltage latch, sub-threshold
integration, sharp firing threshold. Three open: knee-gate
$V_{G2}$ (numerical), autonomous self-reset, and free-running
relaxation oscillation in the few-MHz band reported for the cell.

\item \textbf{Networks.} Seven applications pass pre-registered
gates on independent benchmarks, all on the same calibrated cell,
no per-task device tuning: sparse-ER reservoir on Mackey--Glass
(NRMSE~$0.015$ at $N\!=\!1024$); HDC on UCI~HAR ($\approx 84\%$ at
$D\!=\!8192$); associative-memory binder ($\approx 87\%$ recall,
$32$-pair capacity); predictive coding on NAB ($F1\!=\!0.34$);
$1/f$-noise stochastic RNG (NIST $5/5$); $16$-tap LMS equaliser
($\sim\!170\times$ lower energy/symbol vs.\ fp32); KWS$\to$ECG edge
cascade ($F1\!=\!0.845$, $\approx\!61\%$ energy savings vs.\
always-on).
\end{enumerate}

\textbf{What this proposal funds.}
(i) Close the transient gap: add the distributed body-resistance
network ($10$~k$\Omega$--$1$~M$\Omega$ tunable $R_B$) and the
interface/inversion-layer body-capacitance contribution missing from
junction-only $C_j$. Both fit inside the existing differentiable
port; the work item is parameter identification, not new
infrastructure. Closing the gap promotes the artifact from
DC-faithful to transient-faithful.
(ii) Re-evaluate the seven network demonstrations under realistic
device variation, shared-bulk-rail coupling, and silicon $1/f$ noise
to produce a fabrication-survivable topology recommendation.
(iii) Two characterisation runs on existing silicon
($I_c/I_b$ at saturation; pulsed-$V_d$ / TLP for body
$R_b\!\cdot\!C_b$) convert the present fit-defensible parameter
point into a silicon-calibrated one and pin the body-state surrogate
to measured device behaviour, opening the path to a tape-out-ready
cell-parameter recommendation by M9.

\textbf{Timeline.}
M3 (now, achieved): functional artifact at the numbers above.
M6 (Sep~2026): nine-of-nine dynamic battery; network re-evaluation
under realistic noise.
M9 (Dec~2026): fan-out / shared-bulk-rail validation circuit
($10$--$30$ cells) and tape-out-ready parameter recommendation.
M12 (Mar~2027): thick-oxide neuron-cell variant modelled and
benchmarked, completing the synapse--soma--linear-input cell family.

\textbf{Budget.}
Approximately \textbf{$200$~k\euro\ / $12$~months}.
Bergvall $0.8$~FTE ($140$~k\euro); Luciani $0.3$~FTE optional
($35$~k\euro, Julia/Enzyme cross-validation);
Pazos and Lanza in-kind via KAUST. Travel $9$~k\euro; hardware
contingency $5$~k\euro; open-access $3$~k\euro; indirect $\approx 5\%$.
Compute (HP~Z2 + AMD~Ryzen~AI~Max+~$395$ + Radeon~$8060$S iGPU) is
already in place.

\vspace{2pt}\hrule\vspace{2pt}

{\footnotesize\emph{Defensible.}
(i)~a differentiable, GPU-batched Python port of the cell, validated
against ngspice; (ii)~cell-wide DC $\approx 1.2$~decade on full
$25$-bias, forward$+$backward; (iii)~six of nine dynamic behaviours
passing, with the two open transient targets explicitly named;
(iv)~seven network demonstrations passing pre-registered gates;
(v)~a concrete two-measurement path to a silicon-calibrated
parameter recommendation.
\emph{Not defended.} ``NS-RAM is a better reservoir than software
ESN'' (false on like-for-like benchmarks; the claim is energy, not
algorithmic superiority); any $<1.0$~decade DC headline (forward-only
or bias-subset); a fully autonomous oscillator on the present cell
(open; first item on the M6 list).}

\end{document}
